Benchmark contamination evades standard detection when training data is paraphrased---models trained on rephrased MMLU achieve 85.9\% accuracy while bypassing n-gram filters. We investigate whether contamination leaves behavioral signatures detectable via confidence patterns. We propose the Semantic Saturation Index (SSI), measuring confidence variance across paraphrases of benchmark items, hypothesizing that contaminated models exhibit uniform confidence due to robust semantic representations. Through controlled experiments on Mistral-7B (0--50\% MMLU contamination, K=20 paraphrases), we validate the causal mechanism: contamination creates training exposure (31.1\% accuracy gain), diverse training produces representation invariance (MPS difference 0.065, Cohen's d=0.52), and invariance correlates with confidence uniformity (r=$-$0.517). However, the SSI metric itself fails---achieving AUC=0.506 on real data, essentially chance. The inverse-variance formulation amplifies noise rather than signal. We contribute the first empirical demonstration of the contamination-to-uniformity mechanism chain, document the failure of inverse-variance metrics guiding future work toward entropy-based alternatives, and identify a simulation-reality gap where simulated experiments passed but real inference revealed metric failure.
